% Guia do professor — Capítulo 2: Radiação Solar e Infravermelha
\section{Capítulo 2 --- Radiação Solar e Infravermelha}

\subsection*{Visão geral e ritmo}

O capítulo energético do curso: geometria Sol--Terra, leis de corpo negro,
lei de Beer e balanço radiativo, culminando na temperatura efetiva e no
efeito estufa. \textbf{Ritmo sugerido: 2 aulas de 2\,h.}

\begin{itemize}
  \item \textbf{Aula 1} — motivação (radiação como combustível), declinação,
        elevação solar, duração do dia, fluxo e inverso do quadrado.
        Fechar com o Caso~2 (mapa de insolação — 10\,min).
  \item \textbf{Aula 2} — Planck, Wien, Stefan--Boltzmann, Kirchhoff, Beer,
        balanço na superfície, $T_e = 255$\,K e o efeito estufa.
        Fechar com o Caso~4 (climlab — 15\,min).
\end{itemize}

\subsection*{Objetivos de aprendizagem}

\begin{enumerate}
  \item calcular declinação, elevação solar e duração do dia para qualquer
        latitude e data;
  \item aplicar Planck/Wien/Stefan--Boltzmann e justificar a separação
        onda curta / onda longa;
  \item deduzir e aplicar a lei de Beer (forma diferencial e integrada);
  \item calcular $T_e$ de um planeta e explicar quantitativamente o efeito
        estufa (modelo de 1 e $N$ camadas);
  \item montar o balanço radiativo de superfície com albedo e IV atmosférico.
\end{enumerate}

\subsection*{Pré-requisitos a reativar}

Trigonometria esférica mínima (a eq.\ da elevação pode ser \emph{dada});
integral de $x^3/(e^x-1)$ apenas citada (solução simbólica no notebook de
soluções); lei dos gases não é usada — capítulo independente do Cap.~1
exceto pela estrutura vertical.

\subsection*{Armadilhas conceituais frequentes}

\begin{itemize}
  \item \textbf{``Verão porque a Terra está mais perto do Sol''}: o periélio
        é em janeiro — e o verão do hemisfério norte é em julho. O que manda
        é $\sin\Psi$ e a duração do dia. Vale perguntar antes de explicar.
  \item \textbf{Fator 4 do $T_e$}: disco intercepta, esfera emite. Alunos
        esquecem e obtêm 279\,K em vez de 255\,K (sem albedo) — apontar que
        os dois erros (fator e albedo) quase se cancelam.
  \item \textbf{$T_e$ vs.\ temperatura da superfície}: $T_e$ é a temperatura
        de \emph{emissão para o espaço} (nível médio $\sim$5\,km), não do
        chão.
  \item \textbf{Wien com $T$ em °C}: erro clássico; reforçar kelvins.
  \item \textbf{Beer para radiação difusa}: a lei vale para o feixe direto;
        céu nublado exige tratamento de espalhamento múltiplo (fora do
        escopo).
\end{itemize}

\subsection*{Demonstração de baixo custo}

Termômetro IV de cozinha (ou câmera térmica de celular): medir o ``céu''
limpo (dá $-20$ a $-50\,°$C — o $I\!\downarrow$ baixo do ar seco) e uma
nuvem (bem mais quente). Conecta emissividade, $I\!\downarrow$ e efeito
estufa em 5 minutos.

\subsection*{Problemas sugeridos do livro (seção 2.9)}

Aquecimento: A2 (declinação), A7 (elevação), A12--A14 (Planck/Wien/S-B),
A19 (Beer), A22 (balanço). Aprofundamento: E8 (insolação diária),
E15 (albedo e $F^*$). Síntese: S3 (``e se o eixo da Terra fosse
$45°$?'') — casa com o Caso~2. \emph{Confira a numeração na sua cópia
(v1.02b) antes de publicar.}

\subsection*{Uso do notebook \texttt{cap02\_radiacao.ipynb}}

\begin{center}
\begin{tabular}{lll}
\hline
Caso & Conteúdo & Momento sugerido \\
\hline
1 & Curvas de Planck Sol × Terra; Wien e S--B verificados & Aula 2, após Planck \\
2 & Elevação solar, duração do dia, mapa de insolação lat × dia & fim da Aula 1 \\
3 & Lei de Beer como EDO; atenuação vs.\ ângulo zenital & Aula 2, após Beer \\
4 & Estufa de $N$ camadas × coluna cinza \texttt{climlab} & fim da Aula 2 \\
\hline
\end{tabular}
\end{center}

O Caso~2 rende o gráfico mais bonito do início do curso (mapa de insolação
latitude × dia); o Caso~4 introduz o \texttt{climlab}, que voltará no
Cap.~21.

\subsection*{Exercícios complementares e soluções}

Nas notas: T1--T5 (Wien e S--B a partir de Planck, $T_e$ planetário, estufa
de 1 camada, massa de ar $1/\sin\Psi$) e N1--N4 (células reservadas no
notebook). Soluções completas em \texttt{cap02\_solucoes.ipynb}
(\emph{não distribuir}): T1--T2 verificados simbolicamente com
\texttt{sympy}, N4 compara a fórmula $(N{+}1)^{1/4}$ com o modelo cinza.
Par teoria--numérica: T4$\leftrightarrow$N4.

\subsection*{Conexões com o resto do curso}

$F^*$ alimenta o balanço de calor da superfície (Cap.~3); insolação
diferencial $\to$ circulação geral (Cap.~11); Beer $\to$ satélites
(Cap.~8); climlab e estufa $\to$ clima (Cap.~21).
